% !TEX root = /home/eb/.vscode/extensions/qftrules.programme-colle/tmp/Exercice.tex
%%%%%%%%%%%%%%%%%%%%%%%%%%%%%%%%%%%
\begin{exo}[PC][1][colle]{Calculs de différences de marche}
Considérons une étoile $\rm S$ supposée ponctuelle et située à l'infini hors de l'axe optique. Les points $\rm P$ et $\rm Q$ sont donc en phase. L'observation de cette étoile est réalisée par le dispositif où un prisme d'indice $n$ et d'angle $\alpha$ est utilisé pour rediriger les rayons lumineux provenant de l'astre vers un capteur situé au foyer image $\rm F'$ d'une lentille mince convergente de diamètre $a$. Dans tout l'exercice, on considère que l'indice de l'air est $n_a=1$.
Déterminer en justifiant l'expression des différences de marche suivantes :
\begin{fig}[]
  \centering
  \includegraphics[width=0.7\columnwidth]{prism_chemin.pdf}
  % \capt{}
\end{fig}
\begin{quest}
  \item	
  $\opt{PR}-\opt{QS}$;
  \item 
  $\opt{PT}-\opt{QU}$;
\item $\opt{PF'}-\opt{QF'}$.
\end{quest}
\end{exo}
%%%%%%%%%%%%%%%%%%%%%%%%%%%%%%%%%%%%%%%%%%%%

%%%%%%%%%%%%%%%%%%%%%%%%%%%%%%%%%%%
\begin{exo}[PC][1][colle]{Effet Doppler}
  \Notions{\item Onde monochromatique \task Vibration lumineuse \task Effet Doppler}
Une onde monochromatique se propage dans le vide. La vibration en un point $\mathrm{M}(x, y, z)$ à l'instant $t$, exprimée en unités arbitraires de distance et de temps, s'écrit :
\begin{equation}
  \vib(x, y, z, t)=A \cos (2 t-3 x+4 y+5 z).
\end{equation}
\begin{quest}
  \item Que vaut la pulsation $\omega$ ? Que vaut le vecteur d'onde $\Vec{k}$ ? De quelle nature est la surface d’onde ? Déterminer la célérité $c$ de cette onde.
  % \item Montrer que $\Psi(x, y, z, t)$ satisfait à l'équation d'onde de D'alembert. Calculer la vitesse de propagation.
  \item Un observateur se déplace sur l'axe des abscisses $x$ à la vitesse $v$. Montrer que l'observateur perçoit alors une onde de même vecteur d’onde $\Vec{k}$ mais de pulsation différente $\omega^{\prime}$  à déterminer. Ce phénomène s’appelle effet Doppler. Quelle est sa manifestation dans le cas des ondes acoustiques ?
\end{quest}
\end{exo}
%%%%%%%%%%%%%%%%%%%%%%%%%%%%%%%%%%%

%%%%%%%%%%%%%%%%%%%%%%%%%%%%%%%%%%%
\begin{exo}[PC][1][colle]{Élargissement collisionnel d’une raie spectrale}
  \Notions{\item Source lumineuse \task Élargissement collisionnel}
  On considère une source spectrale constituée d’un vapeur atomique, assimilée à un gaz parfait monoatomique de masse molaire $M$, pression $p$, densité particulaire $n$, et température $T$. Les atomes sont assimilés à des sphères dures identiques de masse $m$ et de rayon $r$. On souhaite étudier l’élargissement d’une raie spectrale de fréquence $\nu_0$ émise par les atomes de ce gaz provoqué par les collisions entre atomes.
  \begin{quest}
    \item Déterminer la vitesse quadratique moyenne $v_q$ des atomes en fonction de la température $T$, de la masse $m$, et de la constante des gaz parfait $R$.
    \item On définit la section efficace de collision $\sigma$ comme la surface que peut occuper un centre atomique pour qu’une collision ait lieu avec un autre atome. Exprimer $\sigma$ en fonction du rayon $r$ et de la densité~$n$.
    \item	On suppose qu’un atome cesse d’émettre un train d’onde lorsqu’il collisionne avec un atome. Déterminer la durée moyenne $\tau_c$ d’un train d’onde, puis montrer que la largeur spectrale en fréquence vaut
    $$
    \Delta \nu \sim \sigma p \sqrt{\frac{3}{M RT}}.
    $$
    \AN~: 
  \end{quest}
\end{exo}
%%%%%%%%%%%%%%%%%%%%%%%%%%%%%%%%%%%

%%%%%%%%%%%%%%%%%%%%%%%%%%%%%%%%%%%%%%%%%%%%%%%%%%%%%%%
\begin{exo}[PC][1][TD]{Relations de Descartes}
  \Notions{\item Onde plane \task Théorème de Malus \task Chemin optique \task Snell-Descartes}
Une onde plane de longueur d'onde dans le vide $\lambda_0$ arrive sous un angle d'incidence $\theta$ sur un dioptre plan séparant deux milieux d'indice $n$ et $n'>n$. On suppose que l'onde transmise est plane également, mais se propage dans la direction $\theta'$. On se place dans le cadre de l'optique scalaire.
%\begin{figure}[h!]
%\centering
\begin{center}
\includegraphics[width=0.4\textwidth]{dioptre_descartes.pdf}
\end{center}
%\end{figure}
\begin{quest}
  \item Rappeler l'expression de la vibration lumineuse $\vib(M,t)$ associée à une onde plane. Exprimer le vecteur d'onde $\Vec{k}_i$ et la pulsation $\omega_i$ dans chaque milieu $i$ en fonction des données du problème. Représenter plusieurs plans d'onde associés aux rayons lumineux, en justifiant votre construction.
  \sol{
    La vibration lumineuse associée à une onde plane est 
    \eqbox{
      \vib(M,t) = A\cos(\omega t - \Vec{k}\cdot\Vec{r} -\varphi),
      }
      avec $\Vec{k} = \frac{2\pi n}{\lambda_0} \Vec{u}$ le vecteur d'onde, et $\omega = \frac{2\pi c}{\lambda_0}$ la pulsation.
      Dans le milieu incident d'indice de réfraction $n$, la vibration s'écrit $\vib_1(M,t) = A\cos(\omega t - \Vec{k}_1\cdot\Vec{r} -\varphi_1)$ avec $\Vec{k}_1 = \frac{2\pi n}{\lambda_0} \Vec{u}_1$. Dans le milieu d'indice de réfraction $n$', la vibration s'écrit $\vib_2(M,t) = A\cos(\omega t - \Vec{k}_2\cdot\Vec{r} -\varphi_2)$ avec $\Vec{k}_2 =  \frac{2\pi n'}{\lambda_0} \Vec{u}_2$ et même $\omega$. Les plans d'onde s'obtiennent avec le théorème de Malus : ce sont des plans orthogonaux aux rayons, donc de vecteurs normaux $\Vec{u}_1$ et $\Vec{u}_2$.
      %Ils sont distants de $\lambda$ ou $\lambda' < \lambda$ car $\theta' < \theta$ et donc $n'>n$.
      }
      
      \item Représenter le point $C$  appartenant au rayon passant par $B$ et tel que $\varphi(C)=\varphi(A)$. Faire de même pour le point~$D$ appartenant au rayon passant par $A$ et tel que $\varphi(D)=\varphi(B)$.
      \sol{
        On utilise le théorème de Malus : les plans d'onde sont orthogonaux aux rayons, donc les points $C$ et $D$ sont les intersections des plans d'onde avec les rayons passant par $B$ et $A$ respectivement.
    \begin{fig}[]
      \centering
      \includegraphics[width=0.5\textwidth]{descartes_phase.pdf}
      \capt{}
    \end{fig}
        }
      \item Retrouver alors la loi de Snell-Descartes de la réfraction à partir de la notion de chemin optique.
      
      \sol{
Comme $\varphi(B) - \varphi(C) = \frac{2\pi}{\lambda_0} (CB)$ et $\varphi(B) - \varphi(C) = \varphi(D) - \varphi(A) = \frac{2\pi}{\lambda_0} (AD)$, il y a égalité de chemin optique $\opt{CB}=\opt{AD}$. Donc par projection $n_1 AB \sin \theta_1 = n_2 AB \sin \theta_2$. Cela est vrai pour toute distance $AB$, et donc
\eqbox{
  n_1 \sin \theta_1 = n_2 \sin \theta_2.
}
}
\end{quest}
\end{exo}
%%%%%%%%%%%%%%%%%%%%%%%%%%%%%%%%%%%%%%%%%%%%%%%%%%%%%%%%%%%%%%%%

% %%%%%%%%%%%%%%%%%%%%%%%%%%%%%%%%%%%
% \begin{exo}[PC][1][appli]{Calcul d'intensité lumineuse}
%   \Notions{\item Intensité lumineuse \task Superposition d'ondes}
%  Pour chacune des vibrations lumineuses monochromatiques suivantes, déterminer l'intensité lumineuse mesurée par un photorécepteur dont le temps de réponse $\tau_r$ est très supérieur à la période propre $T$ de l'onde.

%  \begin{enonce}
%   $\vib(x,t) = A \cos(\omega t)\cos(kx)$.
%  \end{enonce}

%  \begin{corrige}
%   On a
%   \eqbox{
%     I = \frac{1}{2} A^2 + \frac{1}{2} A^2 = A^2.
%   }
%  \end{corrige}

%   \begin{enonce}
%     $\vib(x,t) = A \cos(\omega t - kx) +  B \cos(\omega)$.
%   \end{enonce}
%   \begin{corrige}
%     On a
%     \eqbox{
%       I = \frac{1}{2} A^2 + \frac{1}{2} A^2 - A^2 = 0.
%     }
%   \end{corrige}

%   \begin{enonce}
%     $\vib(x,t) = A \cos(\omega t - kx) + A \cos(\omega t + kx + \frac{\pi}{2})$.
%   \end{enonce}


% \end{exo}
% %%%%%%%%%%%%%%%%%%%%%%%%%%%%%%%%%%%

%%%%%%%%%%%%%%%%%%%%%%%%%%%%%%%%%%%%%%%%%%%%%%%%%%%%%%%%%%%%%%%%
\begin{exo}[PC][2][TD]{Relation de conjugaison}
  \Notions{\item Chemin optique \task Lentille mince}
Une lentille mince convergente de distance focale $f'$ est taillée dans un verre d'indice de réfraction $n$. Son épaisseur centrale vaut $e_0$, et son épaisseur vaut $e(r)$ à une distance $r$ du centre. Les centres de la lentille en entrée et en sortie sont notés O$_1$ et O$_2$ respectivement. On souhaite démontrer la relation de conjugaison de Descartes  à partir de la notion de chemin optique. On place pour commencer une source lumineuse ponctuelle à son foyer $\rm F$, comme indiquée ci-dessous. On se place dans les conditions de Gauss, ce qui suppose que les rayons lumineux font des angles très faibles avec l'axe optique.
%\begin{figure}[h!]
%\centering
%\end{figure}

On considère deux rayons lumineux issus de $\rm F$ et passant par la lentille, l’un par le centre $\rm O$, l’autre par un point $\rm I$.  On note $\rm I_1$ et $\rm I_2$ les points d’incidence et de sortie de ce rayon sur la lentille, $\rm M$ un point en sortie de la lentille du rayon passant par $\rm I$, et $\rm H$ son projeté orthogonal sur l'axe optique, à une distance $d$ de $\rm F$.
\begin{center}
\includegraphics[width=0.5\textwidth]{relation_conjugaison.pdf}
\end{center}
\begin{quest}
  \item De quel type est la surface d'onde en amont de la lentille ? En aval de la lentille ? Justifier.
\item Déterminer l’expression approchée du chemin optique $\opt{FH}$.
\sol{
 On a par additivité du chemin optique,
$$\opt{FH} =  \opt{FO_1} + \opt{O_1O_2} + \opt{O_2H}.$$
 Donc $\opt{FM} \simeq f' + ne_0 + (d-f')$ et
\eqbox{
  \opt{FM} \simeq d + ne_0.
}
}
\item Déterminer l’expression approchée du chemin optique $\opt{FM}$. En comparant les phases en $\rm M$ et en $\rm H$, établir l'expression approchée de l'épaisseur de la lentille $e(r)$ pour $r\ll f'$.
\sol{
  Par additivité du chemin optique, on a
$$\opt{FM} = \opt{FI_1} + \opt{I_1I_2} + \opt{I_2M} \simeq  \sqrt{f'^2+r^2} + ne(r) + (d-f').$$
Or $r \ll f’$ dans les conditions de Gauss, et donc 
$$
\opt{FM} \simeq f' + \frac{r^2}{2f'} + ne(r) + (d-f') \simeq d + ne(r) + \frac{r^2}{2f'}.
$$
De plus, $\rm H$ et $\rm M$ appartiennent au même plan d'onde, qui sont orthogonaux à l’axe optique d’après le théorème de Malus, donc $\opt{FH} = \opt{FM}$. Ainsi $\opt{FM} = \opt{FH}$ d’après la formule du déphasage, et donc
\eqbox{
  e(r)\simeq e_0 - \frac{r^2}{2nf'}.
}
}
\end{quest}
On place maintenant en object réel A en amont de la lentille sur l'axe optique, d'image A'.
% , et on modélise la lentille par une lenti
\begin{fig}[label]
  \centering
  \includegraphics[scale=0.8]{relation_conjugaison_objet.pdf}
  %\capt{}
\end{fig}
\begin{quest}
\item  Justifier que le chemin optique $\opt{AA'}$ ne dépend pas de la position du point $\rm I$. Démontrer la relation de conjugaison de Descartes, qui relie les distances algébriques $\overline{OA}$ et $\overline{OA'}$ à la distance focale $f'$.
\sol{
On a $\opt{AA'} = \opt{AO_1} + \opt{O_1O_2} + \opt{O_2A'}  = -\overline{OA} +  \overline{OA'} + ne_0$ qui est indépendant de $r$, et d'autre part
$\opt{AA'} = \opt{AI_1} + \opt{I_1I_2} + \opt{I_2A'} = \sqrt{OA^2+r^2} + ne + \sqrt{OA'^2+r^2} \simeq OA + \frac{r^2}{2OA} + ne_0 - \frac{r^2}{2f'} + OA' + \frac{r^2}{2OA'}$, soit 
$$
0 = \frac{r^2}{2}\left( \frac{1}{OA} - \frac{1}{f'} + \frac{1}{OA'}\right) \forall r,
$$
et donc 
\eqbox{
  \frac{1}{\overline{OA'}} - \frac{1}{\overline{OA}} = \frac{1}{f'} .
  }
}

\end{quest}
\end{exo}
%%%%%%%%%%%%%%%%%%%%%%%%%%%%%%%%%%%%%%%%%%%%%%%%%%%%%%%%%%%%%%%%

%%%%%%%%%%%%%%%%%%%%%%%%%%%%%%%%%%%%%%%%%%%%%%%%%%%%%%%%%%%%%%%%
\begin{exo}[PC][1]{Pouvoir dispersif d'un réseau et d'un prisme}
\begin{quest}
\item Un réseau de pas $a$ est éclairé sous incidence normale par une lumière polychromatique. En notant $D(\lambda)$ l'angle de déviation de l'onde de longueur d'onde $\lambda$, déterminer son pouvoir dispersif défini par $\frac{\der D}{\der \lambda}$ au voisinage de $\lambda$ en fonction de l'ordre $p$ et du pas $a$.
\item Un prisme d'angle au sommet $A=\SI{60}{\degree}$ est taillé dans un verre flint extra dense dont l'indice suit la loi de Cauchy $n(\lambda)=a+b/\lambda^2$ avec $a=1,683$ et $b=\SI{1.3e-2}{(\micro\meter)^{2}}$. En se plaçant au minimum de déviation, exprimer le pouvoir dispersif du prisme, et faire l'application numérique au voisinage de $\lambda_0=\SI{500}{nm}$. Comparer avec un verre crown pour lequel $a=1,518$ et $b=\SI{4,4e-3}{(\micro\meter)^2}$.
\item Déterminer le nombre de traits par millimètres d'un réseau dont le pouvoir dispersif dans l'ordre 1 est du même ordre que celui du prisme. Conclure.
\end{quest}
\end{exo}
%%%%%%%%%%%%%%%%%%%%%%%%%%%%%%%%%%%%%%%%%%%%%%%%%%%%%%%%%%%%%%%%


%%%%%%%%%%%%%%%%%%%%%%%%%%%%%%%%%%%
\begin{exo}[PC][1][appli]{Déphasage}
  \Notions{\item Déphasage \task Chemin optique}
  \hauteurLargeurCadreReponse		{4mm}{6mm}
  Le schéma ci-contre représente la traversée d'une lame de verre, d'épaisseur non uniforme, par une onde plane monochromatique. Cochez les réponses exactes et justifier votre réponse.
  % \begin{wrapfigure}{r}{0.25\textwidth}
    \begin{center}
      \includegraphics[width=0.2\textwidth]{2023_11_27_e357e55eff796372029bg-08}
    \end{center}
  % \end{wrapfigure}
% \begin{quest}
% \item	
\begin{enonce}
  L'onde en A est en phase avec l'onde au point C.
\end{enonce}
\sol{
  Cette affirmation est
\eqbox{Vraie.}
En effet, le chemin optique est exactement le même pour les points A et C, donc $\varphi(A) = \varphi(C)$.
}

\begin{enonce}
  L'onde en A est en phase avec l'onde en B.
\end{enonce}
\sol{
  Cette affirmation est 
\eqbox{Fausse.}
En effet, le rayon passant par $\rm A$ traverse une épaisseur de lame plus grande que celui passant par $\rm B$. 
}
% \item 

\begin{enonce}
  L'onde en A est en avance de phase par rapport à l'onde au point B.
\end{enonce}
\sol{
  Cette affirmation est 
\eqbox{Fausse.}
En effet, la formule du déphasage indique que la phase accumulée est proportionnelle au chemin optique, et donc à la distance $\ell$ et à l'indice $n$ traversé. On a donc $\vphi(A) > \vphi(B)$. La phase totale de la vibration à l'instant $t$ vaut donc $\omega t - \vphi(A) < \omega t - \vphi(B)$ et $A$ est en avance de phase sur $B$. 
}

% \item 
\begin{enonce}
  L'onde en A est en retard de phase par rapport l'onde au point B.
\end{enonce}
\sol{
  Cette affirmation est
\eqbox{Vraie.}
}
% \item

\sol{
  \udl{Démonstration complète}.

  
  En effet, définissons les points $A_1$, $B_1$ et $C_1$ de l'interface air-verre en entrée, et appartenant aux rayons lumineux passant en $A$, $B$ et $C$ respectivement. Définissons de même les points $A_2$, $B_2$ et $C_2$ de l'interface air-verre en sortie.
  % , et appartenant aux rayons lumineux passant en $A$, $B$ et $C$ respectivement.
  
  L'interface air-verre en entrée est un plan d'onde d'après le théorème de Malus, donc $\vphi(A_1) = \vphi(B_1) = \vphi(C_1)=\vphi_0$. Notons $e$ l'épaisseur de la lame au centre, $\ell>e$ l'épaisseur de la lame en haut et en bas, et $L=A_1A_2=B_1B_2=C_1C_2$ les ditances géométriques, $n$ l'indice du verre. L'air est supposé d'indice égal à $1$.
  n a ensuite par définition du chemin optique $\vphi(A_2) = \vphi(A_1) +\frac{2\pi}{\lambda_0}n\ell$, $\vphi(B_2) = \vphi(B_1) + \frac{2\pi}{\lambda_0}ne$ et $\vphi(C_2) = \vphi(C_1) + \frac{2\pi}{\lambda_0}n\ell$, puis $\vphi(A) = \vphi(A_2) + \frac{2\pi}{\lambda_0}(L-\ell)$, $\vphi(B) = \vphi(B_2) + \frac{2\pi}{\lambda_0}(L-e)$ et $\vphi(C) = \vphi(C_2) + \frac{2\pi}{\lambda_0}(L-\ell)$. On en déduit 
  \eqbox{
  \vphi(A)  = \vphi_0 + \frac{2\pi}{\lambda_0}L + \frac{2\pi}{\lambda_0}(n-1)\ell.
  }
  \eqbox{
  \vphi(C)  = \vphi_0 + \frac{2\pi}{\lambda_0}L + \frac{2\pi}{\lambda_0}(n-1)\ell = \vphi(A)
  }
  \eqbox{
  \vphi(B)  = \vphi_0 + \frac{2\pi}{\lambda_0}L + \frac{2\pi}{\lambda_0}(n-1)e < \vphi(A)=\vphi(C)
  }
}
% \end{quest}
\end{exo}
%%%%%%%%%%%%%%%%%%%%%%%%%%%%%%%%%%%

%%%%%%%%%%%%%%%%%%%%%%%%%%%%%%%%%%%%%%%%%%%%%%%%%%%%%%%%%%%%%%%%%%%%%%%
\begin{exo}[PC]{Transformation d'une surface d'onde: collimateur}
Une onde sphérique issue d'une source ponctuelle $S$ est placée au foyer objet d'une lentille mince $\mathcal{L}$ plan-convexe, de rayon de courbure $R$, d'indice de réfraction $n$ et de distance focale $f$. Le milieu ambiant est l'air avec $n\simeq 1$. En écrivant que l'action de la lentille est de déformer la surface d'onde de façon à la rendre plane, on cherche à obtenir l'expression de la distance focale de la lentille en fonction de son indice $n$ et de son rayon de courbure $R$.
\begin{quest}
\item Exprimer l'épaisseur $e$ de verre traversée, en fonction de la distance $y$ à l'axe de la lentille, de son rayon de courbure $R$ et de son épaisseur maximale sur l'axe optique $e_0$. On supposera que $y\ll R$ (rayon paraxial).
\sol{
$(R-e_0+e)^2 + y^2 = R^2$ donc $R-e_0+e = \sqrt{R^2-y^2} \simeq R - \frac{y^2}{2R}.$
}
\item On considère les points $A$ et $A'$, intersections de la surface de la lentille et de l'axe optique. Tracer les surfaces d'ondes passant par ces deux points. On notera $B$ et $B'$ les points de ces deux surfaces d'onde appartenant au rayon coupant la lentille à la distance $y$ de son axe.
\sol{

}
\item Exprimer indépendamment les chemins optiques $\opt{AA'}$ et $\opt{BB'}$ en fonction de $n,e_0,y,f$ et $R$.
\sol{
$\opt{AA'} = ne_0$
}
\item Déduire une relation entre ces chemins optiques, puis déterminer à l'aide de la question précédente $f$ en fonction de $n$ et $R$.
\end{quest}
\end{exo}
%%%%%%%%%%%%%%%%%%%%%%%%%%%%%%%%%%%%%%%%%%%%%%%%%%%%%%%%%%%%%%%%%%%%%%%

%%%%%%%%%%%%%%%%%%%%%%%%%%%%%%%%%%%
\begin{exo}[PC][2][TD]{Approximation paraxiale d'une onde sphérique}
  \Notions{\item Onde sphérique \task Onde plane \task Chemin optique}
  Cherchons à établir un critère quantitatif pour assimiler localement une onde sphérique à une onde plane, en se plaçant dans les conditions de Gauss (éloignement et inclinaison faibles par rapport à l'axe optique). On peut pour cela estimer l'écart amplitude et en déphasage entre l'onde sphérique exacte et l'onde plane approchée.


  %On considère une onde incidente sphérique, d'amplitude complexe $\underline{a}(r)=\frac{A}{r}\,\e^{-i\vv{k}\cdot\vv{r}}$.

  %
   \begin{fig}
    \centering
    \includegraphics[width=0.5\textwidth]{approx_paraxiale.pdf}
    % \capt{Approximation paraxiale d'une onde sphérique.}
  \end{fig}
  %
  Soit $\rm M$ un point quelconque situé à une distance $z$ de la source le long de l'axe optique, et à une hauteur~$h$ par rapport à l'axe optique, $\rm C$ le projeté orthogonal de $\rm M$ sur l'axe optique, et $\rm H$ le point de la droite $(SM)$ situé à une distance $z$ de $\rm S$.
  
  \begin{quest}
    \item	
    En utilisant le théorème de Pythagore et une approximation linéaire ($h\ll z$ d'après les conditions de Gauss), déterminer une expression approchée du déphasage $\Delta\varphi_{MC}=\varphi(M)-\varphi(C)$.
    \sol{
      Comme $H$ et $C$ appartiennent à la même surface d'onde, la vibration lumineuse a même phase en ces points, et $\opt{HM}=\opt{HC}$. On se place dans l'air ($n=1$) et on cherche à évaluer la différence de chemin optique $SM-SC$. 
      \begin{equation}
       SM-SC=\sqrt{z^2+h^2}-z=z\left[\sqrt{1+\left(\frac{h}{z}\right)^2}-1\right]\simeq\frac{h^2}{2z}.
       \end{equation}
      Le déphasage entre $M$ et $C$ (ou $H$) vaut,
      \eqbox{
        \Delta\varphi_{MC} = \varphi_M-\varphi_C=\frac{2\pi}{\lambda_0}(SM-SC)\simeq \frac{\pi h}{\lambda_0}\left(\frac{h}{z}\right).
      }
    }
    \item Déterminer de même une expression approchée de l'amplitude de l'onde sphérique en $M$.
    \sol{
      Au premier ordre en $\frac{h}{z}\ll 1$, l'amplitude de l'onde sphérique peut se simplifier en
      \eqbox{
        \frac{1}{SM}=\frac{1}{z\sqrt{1+\frac{h^2}{z^2}}}\simeq \dfrac{1}{z}\left(1-\frac{h^2}{2z^2}\right) \simeq  \dfrac{1}{z}.
      }
    } 
    \item En déduire une expression de la vibration lumineuse de l'onde sphérique en $M$ dans l'approximation paraxiale. À quelle condition cette onde peut-elle être considérée comme plane ? Ordre de grandeur pour diaphragme décentré d'une $h=\SI{1}{cm}$ au dessus de l'axe optique éclairée en lumière verte ($\lambda=\SI{500}{\nano\meter}$).
    \sol{
    
      Dans les conditions de Gauss, on peut ainsi exprimer l'onde sphérique comme
      \begin{equation}\label{eq:parax}
      \vib(z,h,t)\simeq \frac{a}{z}\cos(\omega t-\varphi_M) = \frac{a}{z}\cos(\omega t-\varphi_C -\Delta\varphi_{MC}) \simeq \frac{a}{z}\cos\!\left(\omega t-kz -\frac{\pi h^2}{\lambda_0 z}\right).
      \end{equation}
      L'\Eq{parax} est l'expression d'une onde sphérique dans l'approximation paraxiale. Cette onde peut être supposé  plane localement (c'est à dire en un point de distance $z$ fixé) si
      \begin{equation}
      \frac{\pi h^2}{\lambda_0 z}\ll 1\quad\Rightarrow\quad \boxed{z \gg \displaystyle\frac{\pi h^2}{\lambda_0}}.
      \end{equation}
      Ordre de grandeur : si le montage est mal aligné, de sorte qu'un élément du montage est décentré d'une hauteur $h=\SI{1}{cm}$ au dessus de l'axe optique, il faudrait se placer à $z\gg\SI{600}{m}$ pour pouvoir considérer l'onde comme plane ($\lambda=\SI{500}{nm}$) ! C'est évidemment irréalisable, c'est pourquoi l'on utilise des lentilles pour « construire » des ondes quasi planes !
    }
    \end{quest}

\end{exo}
%%%%%%%%%%%%%%%%%%%%%%%%%%%%%%%%%%%
